% =============================================================================
% Appendix M: Audited Theorems and Gap Register
% ORIUS/DC3S Battery-Safety Thesis
% =============================================================================
\chapter{Audited Theorems and Gap Register}
\label{app:verified-theorems-gap-audit}
\begingroup
\footnotesize
\sloppy

This appendix is the human-readable companion to the machine-generated
active theorem audit in \texttt{reports/publication/active\_theorem\_audit.*}.
It does \emph{not} treat the live theorem inventory as one flat ladder.
Instead it mirrors the promotion-gated split emitted by
\texttt{theorem\_promotion\_matrix.*}:
\begin{itemize}
  \item \textbf{Flagship theorem}: T1, T2, T3a, T4, T5, T6, T7, T8, T9, T10,
  T11, \texttt{T11\_Byzantine}, \texttt{T\_stale\_decay}, and
  \texttt{T\_trajectory\_PAC}.
  \item \textbf{Flagship corollary}: T3b.
  \item \textbf{Flagship lemma}: L1--L4.
  \item \textbf{Scoped flagship theorem}: \texttt{T\_minimax}.  The sensor
  converse is promoted as \texttt{T\_sensor\_converse}.
\end{itemize}
This appendix audits the full live surface.  Historical drafts remain visible
only as archived text and do not inflate the promoted theorem count.

The certificate half-life remains a corollary, and safe-budget monotonicity
remains a proposition. The tracked integrated theorem surface register still
audits the older hard-gated traceability surface, but it is a
\emph{traceability gate}, not a proof-soundness gate.

\section{Theorem Statements}

\subsection{Theorem 1: OASG Existence}
\label{sec:thm-oasg}

\textbf{Statement.} There exist scenarios where an action appears safe in observed state but is unsafe in true physical state. Formally: $\exists$ episodes with degraded telemetry such that $\text{observed-safe}(a \mid x_{\text{obs}}) \neq \text{true-safe}(a \mid x_{\text{true}})$.

\textbf{Assumptions.} A2 (bounded telemetry error), A4 (known dynamics). Fault injection produces non-trivial $x_{\text{obs}} \neq x_{\text{true}}$ separation.

\textbf{Proof status.} Structural: code explicitly separates observed and true state; empirical: violation rates under dropout/drift exceed nominal.

\textbf{Code anchor.} \texttt{src/orius/cpsbench\_iot/}: \texttt{scenarios.py} (fault generation, $x_{\text{obs}}$, $x_{\text{true}}$, \texttt{event\_log}); \texttt{plant.py} (unclamped truth SOC in \texttt{BatteryPlant.step()}); \texttt{runner.py} (true-state violation computation).

\subsection{Theorem 2: One-Step Safety Preservation}
\label{sec:thm-safety-preservation}

\textbf{Statement.} If the repaired action lies inside the tightened safe set
and the current state lies inside the inflated certificate, then the next
battery state remains safe.

\textbf{Assumptions.} A1--A5, A7.

\textbf{Proof status.} Proved as a conditional one-step theorem plus
deterministic runtime guarantee checks.  Coverage frequency is supplied
separately and is not bundled into T2.

\textbf{Code anchor.} \texttt{src/orius/dc3s/guarantee\_checks.py} and
\texttt{src/orius/dc3s/shield.py}.

\subsection{Theorem 3a: ORIUS Core Envelope Derivation}
\label{sec:thm-core-bound}

\textbf{Statement.} The active defended T3a surface is the per-step envelope
\[
  \Prob[Z_t = 1 \mid \mathcal{H}_t] \le \alpha(1-w_t),
\]
under the explicit battery risk-budget contract and the narrowed
reliability-score interpretation.

\textbf{Assumptions.} A1, A2, A4, A5, A6, A7, A9, plus the theorem-local
battery calibration bridge recorded in the active theorem audit.

\textbf{Proof status.} Paper-rigorous flagship theorem with explicit runtime
contract reporting.  The code now emits a theorem-contract summary that checks
the executable envelope algebra and keeps the calibration bridge explicit as a
declared theorem-local contract rather than hiding it inside generic conformal
language.

\textbf{Code anchor.} \texttt{src/orius/universal\_theory/risk\_bounds.py}
(\texttt{build\_t3a\_contract\_summary()}), plus the backward-compatible
wrapper \texttt{src/orius/dc3s/coverage\_theorem.py}.

\subsection{Corollary 3b: ORIUS Core Aggregation Corollary}
\label{sec:thm-core-bound-aggregation}

\textbf{Statement.} Aggregating the predictable per-step budget from T3a gives
the episode envelope
\[
  \mathbb{E}[V] \le \alpha(1-\bar{w})T.
\]

\textbf{Assumptions.} The assumptions of T3a, plus linearity of expectation.

\textbf{Proof status.} Supporting defended corollary.  Once T3a is fixed, the
aggregation is algebraic bookkeeping rather than a new theorem burden.

\textbf{Code anchor.} \texttt{src/orius/universal\_theory/risk\_bounds.py}
(\texttt{compute\_episode\_risk\_bound()}).

\subsection{Theorem 4: Observation Necessity / No Free Safety}
\label{sec:thm-no-free}

\textbf{Statement.} Within the fixed-margin, quality-ignorant controller
class, there exists an admissible degraded-observation sequence that produces
an observation-action safety gap and hence a true-state violation.

\textbf{Assumptions.} A3 (feasible safe repair exists), A8 (admissible fallback policy). Quality-ignorant baselines lack reliability-aware inflation.

\textbf{Proof status.} Constructive witness theorem on the battery reference
surface, plus empirical ablation support.

\textbf{Code anchor.} Supporting evidence in \texttt{baselines.py},
\texttt{scenarios.py}, and \texttt{runner.py}.

\subsection{Definition 5: Certificate Validity Horizon}
\label{sec:thm-cert-horizon}

\textbf{Statement.} T5 is kept as the constructive maximal forward-tube
containment definition: the validity horizon is the largest integer horizon
whose forward tube remains inside the battery SOC bounds.

\textbf{Assumptions.} A1, A2, A4--A7.

\textbf{Proof status.} Explicit definition, not part of the defended theorem
count.  The helper computes the maximal containment horizon directly and no
probabilistic future-step theorem is claimed at this surface.

\textbf{Code anchor.} \texttt{src/orius/dc3s/temporal\_theorems.py}
(\texttt{forward\_tube()}, \texttt{certificate\_validity\_horizon()});
registered as T5 in \texttt{THEOREM\_REGISTER}
(\texttt{src/orius/dc3s/theoretical\_guarantees.py}).
Runtime expiry logic consumes this quantity, but the active theorem audit does
not count T5 as a defended theorem.

\subsection{Theorem 6: Certificate Expiration Bound}
\label{sec:thm-cert-expiration}

\textbf{Statement.} The battery certificate horizon satisfies the canonical
delta-aware lower bound
\[
  \tau_t \ge
  \left\lfloor
    \frac{\delta_{\mathrm{bnd},t}^2}
    {2\sigma_d^2\log(2/\delta)}
  \right\rfloor.
\]

\textbf{Assumptions.} A4--A7.

\textbf{Proof status.} Flagship defended closed-form theorem with
machine-checkable helper linkage.  The theorem-facing API requires
\texttt{delta} explicitly, records the first-passage side conditions, and
marks legacy no-\texttt{delta} semantics as non-defended.

\textbf{Code anchor.} \texttt{src/orius/universal\_theory/battery\_instantiation.py}
(\texttt{certificate\_expiration\_bound()}), which now emits theorem metadata
and a fail-closed T6 contract summary.

\subsection{Theorem 7: Feasible Fallback Existence}
\label{sec:thm-fallback}

\textbf{Statement.} The defended T7 surface is piecewise: zero dispatch
preserves safety on the interior, a boundary-aware safe-landing action
preserves safety near the boundary, and the runtime fails closed when neither
case can be certified.

\textbf{Assumptions.} A1, A4, A8.

\textbf{Proof status.} Flagship defended battery fallback theorem with an
explicit piecewise executable witness.  The promoted surface is battery-only
and does not claim cross-domain transfer on its own.

\textbf{Code anchor.} \texttt{src/orius/dc3s/temporal\_theorems.py}
(\texttt{zero\_dispatch\_fallback()}, \texttt{certify\_fallback\_existence()}).

\subsection{Theorem 8: Graceful Degradation Dominance}
\label{sec:thm-graceful-dominance}

\textbf{Statement.} The defended T8 surface is the paired sequence-level
comparison: for graceful and uncontrolled violation sequences produced under
the same admissible fault trace,
\[
  z_t^{\mathrm{graceful}} \le z_t^{\mathrm{uncontrolled}} \;\; \forall t
  \quad\Longrightarrow\quad
  V^{\mathrm{graceful}} \le V^{\mathrm{uncontrolled}}.
\]
Strict improvement is claimed only when the graceful sequence is strictly
smaller on at least one step.

\textbf{Assumptions.} A1--A5, A8.

\textbf{Proof status.} Supporting defended comparison theorem at the
sequence-comparison level only.  The stronger absorbing-landing interpretation
remains supporting evidence rather than a promoted theorem.

\textbf{Code anchor.} \texttt{src/orius/universal\_theory/battery\_instantiation.py}
(\texttt{evaluate\_graceful\_degradation\_dominance()}), with
\texttt{safe\_landing} kept as supporting implementation evidence only.

\subsection{Theorem 9: Universal Impossibility}
\label{sec:thm-universal-impossibility}

\textbf{Statement.} For any quality-ignorant controller~$\pi$ and any
$\varphi$-mixing fault process with mean degradation rate~$d > 0$,
$\mathbb{E}[V_T(\pi)] \ge c\,d\,T$ for some constant $c > 0$.

\textbf{Assumptions.} A4 (boundary reachability), A6$'$ ($\varphi$-mixing
fault process).

\textbf{Proof status.} Structural impossibility theorem under explicit
$\varphi$-mixing and boundary-reachability assumptions; battery T4 is the
reference witness for the per-window violation probability and the witness
constant remains theorem-local.

\textbf{Code anchor.} No dedicated module; the battery
\texttt{cpsbench\_iot/runner.py} provides the reference witness;
the proof is summarized in Section~\ref{sec:thm-universal-impossibility}.

\subsection{Theorem 10: Boundary-Indistinguishability Lower Bound}
\label{sec:thm-info-lower-bound}

\textbf{Statement.} Under the explicit boundary-mass and
indistinguishability hypotheses of T10, every controller~$\pi$ satisfies the
scoped lower bound
$\mathbb{E}[V_T(\pi)] \ge \tfrac{1}{4}\,\alpha(1-\bar{w})T$.

\textbf{Assumptions.} A1, A2, A5$'$ (calibration regularity), conformal
coverage condition.

\textbf{Proof status.} Narrowed stylized lower-bound surface under extra
calibration/indistinguishability assumptions.  The active audit records that
the one-sided Le~Cam step remains open, so the code witness is treated as a
scoped lower-bound calculator rather than a sharp frontier theorem.

\textbf{Code anchor.} \texttt{src/orius/universal/reliability\_safety\_frontier.py}
(\texttt{compute\_frontier}, \texttt{is\_dc3s\_optimal}); 25 unit tests
in \texttt{tests/test\_reliability\_frontier.py}.

\subsection{Theorem 11: Typed Structural Transfer}
\label{sec:thm-universality-completeness}

\textbf{Statement.} Any CPS domain instantiation satisfying the structural
transfer obligations of T11---coverage of the observation-consistent state
set, soundness of the tightened safe-action set, repair membership, and
safe fallback admissibility---inherits the T2--T3 safety-guarantee pattern
without new battery-specific algebra.

\textbf{Assumptions.} The universal adapter contract and the four structural
obligations stated in Section~\ref{sec:thm-universality-completeness}.

\textbf{Proof status.} Forward one-step transfer theorem.  The runtime now
builds the four T11 obligation checks directly from typed artifacts, while the
converse is tracked separately as a structural failure-mode proposition rather
than as part of the defended T11 statement.

\textbf{Code anchor.} \texttt{src/orius/universal\_theory/contracts.py} and
\texttt{src/orius/universal\_theory/kernel.py}, with a compact reference
harness in \texttt{src/orius/universal/contract.py}; focused tests in
\texttt{tests/test\_universal\_contract.py},
\texttt{tests/test\_unification.py}, and
\texttt{tests/test\_reliability\_frontier.py}, plus
\texttt{tests/test\_universal\_theory.py}.

\section{Gap Audit Table}

Table~\ref{tab:theorem-gap-audit} summarises the promoted active T1--T11
audit surface.  The publication-ready defended register is the stricter
promotion matrix, where every active theorem-like row has proof, code, tests,
artifacts, hashes, and a claim boundary.

\begin{table}[htbp]
\centering
\caption{Promoted T1--T11 theorem audit summary after promotion-gate closure}
\label{tab:theorem-gap-audit}
\small
\resizebox{\textwidth}{!}{%
\begin{tabular}{@{}llllp{3.2cm}@{}}
\toprule
\textbf{Theorem} & \textbf{Proof status} & \textbf{Empirical validation} & \textbf{Code matches statement} & \textbf{Notes} \\
\midrule
T1 OASG Existence & Structural & Yes & Yes & Observed-vs-true split is explicit in benchmark code and locked evidence \\
T2 One-Step Safety Preservation & Yes (conditional one-step) & Yes & Yes & Shield + guarantee checks + unit tests anchor the theorem; coverage frequency handled separately \\
T3a ORIUS Core Envelope Derivation & Paper-rigorous & Yes & Yes & Runtime theorem-contract summary now separates executable envelope algebra from the declared calibration bridge \\
T3b ORIUS Core Aggregation Corollary & Corollary only & Yes & Yes & Supporting aggregation corollary derived from T3a by linearity of expectation \\
T4 Observation Necessity / No Free Safety & Constructive witness & Yes & Yes & Battery proof gives an explicit degraded-observation sequence and the promotion card states the mandatory-release scope \\
T5 Finite-Horizon Certificate Validity & Flagship theorem & Yes & Yes & Reachable-tube containment, invalidating events, horizon artifacts, and result card close the promotion gate \\
T6 Certificate Expiration Bound & V2-linked flagship theorem & Yes & Yes & Canonical helper is delta-aware, closed-form, and records explicit first-passage side conditions \\
T7 Feasible Fallback Existence & Piecewise hold-or-safe-landing theorem & Yes & Yes & Runtime certifies interior safe hold, boundary recovery by safe landing, and fail-closed infeasibility \\
T8 Graceful Degradation Dominance with Useful Work & Flagship theorem & Yes & Yes & Paired-trace safety dominance plus useful-work retention is backed by policy comparison artifacts \\
\midrule
T9 Universal Impossibility & Structural / scoped witness constant & Yes & Partial & $\varphi$-mixing + Azuma windowing remain explicit assumptions; battery T4 is the reference witness; Section~\ref{sec:thm-universal-impossibility} \\
T10 Boundary-Indistinguishability Lower Bound & Flagship theorem & Yes & Yes & Two-state TV lower bound is scoped to disjoint safe sets and finite boundary pairs; no global frontier claim \\
T11 Typed Structural Transfer & Forward-only one-step transfer & Yes & Yes & Typed runtime artifacts now carry four-obligation summaries; the converse and mini-harness remain separate supporting surfaces \\
\bottomrule
\end{tabular}%
}
\end{table}

\begin{table}[htbp]
\centering
\caption{Audited extension laws and depth-theorem surfaces}
\label{tab:extension-theorem-gap-audit}
\small
\resizebox{\textwidth}{!}{%
\begin{tabular}{@{}llllp{3.1cm}@{}}
\toprule
\textbf{ID} & \textbf{Status} & \textbf{Code surface} & \textbf{Audit posture} & \textbf{Notes} \\
\midrule
L1 & Flagship lemma & \texttt{runtime\_laws.py} & Promotion-gated & Reliability-margin inflation is monotone in the stated runtime law \\
L2 & Flagship lemma & \texttt{runtime\_laws.py} & Promotion-gated & Safe-set antitonicity is stated as a set-containment lemma \\
L3 & Flagship lemma & \texttt{runtime\_laws.py} & Promotion-gated & Unsafe candidate actions must repair, fallback, or deny release \\
L4 & Flagship lemma & \texttt{runtime\_laws.py} & Promotion-gated & Observation-ambiguity sandwich is a scoped lower/upper law \\
T11\_Byzantine & Flagship theorem & \texttt{quality.py} & Promotion-gated & Trimmed-mean reliability aggregation is valid for $b<n/2$ Byzantine channels \\
T\_stale\_decay & Flagship theorem & \texttt{stale\_decay.py} & Promotion-gated & Stale-hold uncertainty growth follows from bounded latent drift \\
T\_minimax & Scoped flagship theorem & \texttt{minimax\_boundary.py} & Promotion-gated & Finite ambiguity-class lower bound only; no unrestricted minimax optimality claim \\
T\_sensor\_converse & Flagship theorem & \texttt{sensor\_necessity.py} & Promotion-gated & Sensor necessity holds under adapter semantics and disjoint safe sets \\
T\_trajectory\_PAC & Flagship theorem & \texttt{risk\_bounds.py} & Promotion-gated & Bonferroni finite-horizon certificate; no separate Ville-strengthened claim \\
\bottomrule
\end{tabular}%
}
\end{table}

\section{Remaining Claim-Boundary Gaps}

The audit above separates proof completeness from promotion completeness.
The remaining gaps that still bound the dissertation's final claim are
therefore explicit:
\begin{enumerate}
  \item \textbf{Autonomous vehicles.} The current AV row is no longer claimed
    as full autonomy closure. It is promoted only as a bounded nuPlan replay
    row under TTC and predictive-entry-barrier semantics. Remaining AV work is
    closed-loop road deployment evidence, not a contradiction of T1--T11.
  \item \textbf{Future mobility and flight-style rows.} Earlier drafts named
    additional exploratory rows, but the current monograph does not promote
    them. They remain planning surfaces until they supply defended replay,
    runtime, governance, and post-repair evidence under the same gate as the
    active Battery + AV + Healthcare lane.
\end{enumerate}

These are the remaining thesis-level gaps after the monograph theorem pass.
They do not weaken the battery proof witness; they bound how far the transfer
claim is currently allowed to extend.

\section{Locked Empirical Evidence}

Battery theorem verification relies on a mixture of code anchors, tests, and
locked outputs in \texttt{reports/publication/}:
\begin{itemize}
\item \texttt{dc3s\_main\_table\_ci.csv}: violation rates, PICP, scenario breakdown
\item \texttt{reliability\_group\_coverage\_phase3.csv}: coverage by reliability bin
\item \texttt{blackout\_study.csv}: certificate-safe-hold evidence for the temporal extension
\item \texttt{graceful\_degradation\_trace.csv}: behavioral evidence for the landing controller
\item \texttt{claim\_evidence\_matrix.csv}: claim-to-evidence linkage
\end{itemize}

The theorem-to-evidence map is maintained in \texttt{orius-plan/theorem\_to\_evidence\_map.md}.

\section{Integrated Release-Gate Status}

The flagship T1--T8 ladder above remains the governance-critical proof spine,
and the broader defended theorem program now extends through T11.  The
generated summary below reports the older hard-gated release status for the 10
source-backed theorem claims in the precursor plus reference ladder together
with the 8 Appendix~C theorem restatements.  The larger T1--T11 program is
kept aligned by the theorem-surface validator and the appendix proof crosswalk
rather than by that legacy 18-row gate alone.

\IfFileExists{reports/publication/integrated_theorem_gate_summary.tex}{%
  \begin{table}[htbp]
\centering
\caption{Integrated 18-row hard-gated theorem summary. The unique theorem set
contains the two supporting Chapter~4 theorem statements and the
flagship T1--T8 ladder; the appendix class checks the eight theorem
restatements in Appendix~C against that unique surface.}
\label{tab:integrated-theorem-gate-summary}
\begin{tabular}{lrr}
\toprule
Class & Passed & Total \\
\midrule
Unique theorem claims & 10 & 10 \\
Appendix theorem restatements & 8 & 8 \\
\midrule
\textbf{All hard-gated theorem rows} & \textbf{18} & \textbf{18} \\
\bottomrule
\end{tabular}
\end{table}
%
}{%
  \begin{table}[htbp]
\centering
\caption{Integrated 18-row hard-gated theorem summary. The unique theorem set
contains the two supporting Chapter~4 theorem statements and the
flagship T1--T8 ladder; the appendix class checks the eight theorem
restatements in Appendix~C against that unique surface.}
\label{tab:integrated-theorem-gate-summary}
\begin{tabular}{lrr}
\toprule
Class & Passed & Total \\
\midrule
Unique theorem claims & 10 & 10 \\
Appendix theorem restatements & 8 & 8 \\
\midrule
\textbf{All hard-gated theorem rows} & \textbf{18} & \textbf{18} \\
\bottomrule
\end{tabular}
\end{table}
%
}

\IfFileExists{reports/publication/integrated_theorem_gate_matrix.tex}{%
  \begin{small}
\begin{longtable}{@{}p{1.3cm}p{2.3cm}p{4.1cm}p{1.4cm}p{1.4cm}p{1.4cm}p{1.4cm}p{1.8cm}@{}}
\caption{Integrated theorem release-gate matrix.}
\label{tab:integrated-theorem-gate-matrix}\\
\toprule
Key & Class & Title & Source & Code & Artifacts & Mapping & Status \\
\midrule
\endfirsthead
\multicolumn{8}{c}{\tablename\ \thetable{} --- continued}\\
\toprule
Key & Class & Title & Source & Code & Artifacts & Mapping & Status \\
\midrule
\endhead
\midrule
\multicolumn{8}{r}{\textit{Continued on next page}}\\
\endfoot
\bottomrule
\endlastfoot
S1 & unique\_theorem & Existence of the illusion under dropout & yes & yes & yes & yes & pass \\
S2 & unique\_theorem & DC3S Feasibility Guarantee & yes & yes & yes & yes & pass \\
T1 & unique\_theorem & OASG Existence & yes & yes & yes & yes & pass \\
T2 & unique\_theorem & Safety Preservation & yes & yes & yes & yes & pass \\
T3 & unique\_theorem & ORIUS Core Bound & yes & yes & yes & yes & pass \\
T4 & unique\_theorem & No Free Safety & yes & yes & yes & yes & pass \\
T5 & unique\_theorem & Certificate validity horizon & yes & yes & yes & yes & pass \\
T6 & unique\_theorem & Certificate expiration bound & yes & yes & yes & yes & pass \\
T7 & unique\_theorem & Feasible fallback existence & yes & yes & yes & yes & pass \\
T8 & unique\_theorem & Graceful degradation dominance & yes & yes & yes & yes & pass \\
C.1 & appendix\_restatement & C.1 Battery-Domain OASG Existence & yes & yes & yes & yes & pass \\
C.2 & appendix\_restatement & C.2 One-Step Safety Preservation & yes & yes & yes & yes & pass \\
C.3 & appendix\_restatement & C.3 Battery ORIUS Core Bound & yes & yes & yes & yes & pass \\
C.4 & appendix\_restatement & C.4 No Free Safety & yes & yes & yes & yes & pass \\
C.5 & appendix\_restatement & C.5 Certificate Validity Horizon & yes & yes & yes & yes & pass \\
C.6 & appendix\_restatement & C.6 Certificate Expiration Bound & yes & yes & yes & yes & pass \\
C.8 & appendix\_restatement & C.8 Feasible Fallback Existence & yes & yes & yes & yes & pass \\
C.9 & appendix\_restatement & C.9 Graceful Degradation Dominance & yes & yes & yes & yes & pass \\
\end{longtable}
\end{small}
%
}{%
  \begin{small}
\begin{longtable}{@{}p{1.3cm}p{2.3cm}p{4.1cm}p{1.4cm}p{1.4cm}p{1.4cm}p{1.4cm}p{1.8cm}@{}}
\caption{Integrated theorem release-gate matrix.}
\label{tab:integrated-theorem-gate-matrix}\\
\toprule
Key & Class & Title & Source & Code & Artifacts & Mapping & Status \\
\midrule
\endfirsthead
\multicolumn{8}{c}{\tablename\ \thetable{} --- continued}\\
\toprule
Key & Class & Title & Source & Code & Artifacts & Mapping & Status \\
\midrule
\endhead
\midrule
\multicolumn{8}{r}{\textit{Continued on next page}}\\
\endfoot
\bottomrule
\endlastfoot
S1 & unique\_theorem & Existence of the illusion under dropout & yes & yes & yes & yes & pass \\
S2 & unique\_theorem & DC3S Feasibility Guarantee & yes & yes & yes & yes & pass \\
T1 & unique\_theorem & OASG Existence & yes & yes & yes & yes & pass \\
T2 & unique\_theorem & Safety Preservation & yes & yes & yes & yes & pass \\
T3 & unique\_theorem & ORIUS Core Bound & yes & yes & yes & yes & pass \\
T4 & unique\_theorem & No Free Safety & yes & yes & yes & yes & pass \\
T5 & unique\_theorem & Certificate validity horizon & yes & yes & yes & yes & pass \\
T6 & unique\_theorem & Certificate expiration bound & yes & yes & yes & yes & pass \\
T7 & unique\_theorem & Feasible fallback existence & yes & yes & yes & yes & pass \\
T8 & unique\_theorem & Graceful degradation dominance & yes & yes & yes & yes & pass \\
C.1 & appendix\_restatement & C.1 Battery-Domain OASG Existence & yes & yes & yes & yes & pass \\
C.2 & appendix\_restatement & C.2 One-Step Safety Preservation & yes & yes & yes & yes & pass \\
C.3 & appendix\_restatement & C.3 Battery ORIUS Core Bound & yes & yes & yes & yes & pass \\
C.4 & appendix\_restatement & C.4 No Free Safety & yes & yes & yes & yes & pass \\
C.5 & appendix\_restatement & C.5 Certificate Validity Horizon & yes & yes & yes & yes & pass \\
C.6 & appendix\_restatement & C.6 Certificate Expiration Bound & yes & yes & yes & yes & pass \\
C.8 & appendix\_restatement & C.8 Feasible Fallback Existence & yes & yes & yes & yes & pass \\
C.9 & appendix\_restatement & C.9 Graceful Degradation Dominance & yes & yes & yes & yes & pass \\
\end{longtable}
\end{small}
%
}
\endgroup

% Theorem-upgrade synchronization note (2026-05-07): see chapters/ch5_theory.tex and appendices/app_c_proofs.tex for completed theorem spine and proof text.
